\chapter{Algorithms}
\label{chapter:algorithms}

The previous chapter derived policy-gradient identities for perturbed mean-field control problems. This chapter describes how these identities are turned into practical algorithms. The emphasis is algorithmic rather than benchmark-specific: choices such as horizons, perturbation grids, validation initial laws, and simulator budgets are deferred to Chapter~\ref{chapter:experiments}.

All methods follow the same outer loop. At iteration $m$, a Monte Carlo estimator $\widehat G_m(\theta_m)$ of the relevant policy gradient is computed, the policy parameter is updated with Adam \cite{kingma2014adam}, and diagnostic quantities such as the objective estimate, gradient norm, and validation objective are recorded. Most benchmarks are written in reward-maximization form; the linear--quadratic benchmark is implemented as a cost minimization and reported with the corresponding sign convention in Chapter~\ref{chapter:experiments}. When a loss-minimization interface is used for a reward objective, the optimizer is applied to the negative estimated gradient.

\section{Baseline policy-gradient estimators}

\paragraph{REINFORCE.}
The first baseline is the standard likelihood-ratio estimator. Given $B$ trajectories generated by the current policy, let $\Rc_\theta^{(b)}$ be the return of trajectory $b$ and let
\begin{align}
    \Sc_\mathrm{pol}^{\theta,(b)}
    =
    \sum_{t=0}^{T-1}
    \nabla_\theta
    \log
    \pi_t^\theta
    \left(
        \alpha_t^{\theta,(b)}
        \mid
        X_t^{\theta,(b)},\mu_t
    \right)
\end{align}
be its policy score. With a deterministic or batch-mean baseline $b_0$, REINFORCE uses
\begin{align}
    \widehat G_{\mathrm{RF}}(\theta)
    =
    \frac1B
    \sum_{b=1}^B
    \left(
        \Rc_\theta^{(b)}-b_0
    \right)
    \Sc_\mathrm{pol}^{\theta,(b)}.
\end{align}
This estimator differentiates the sampled policy but treats the supplied population argument as fixed. In a mean-field control problem it therefore misses the derivative of the population flow with respect to $\theta$.

\paragraph{MF-REINFORCE.}
The second baseline is the finite-state MF-REINFORCE estimator of \cite{meunier2026model}. It perturbs the population law through a logit transformation and uses likelihood-ratio identities to estimate the missing population-flow derivative. This provides a model-free reference estimator for discrete mean-field control, but its auxiliary sampling cost grows quadratically with the time horizon because separate perturbations are used across intermediate time indices.

\section{Transport policy-gradient}

The transport estimator is the practical version of the estimator constructed in Chapter~\ref{chapter:methodology}. It uses two perturbation scales. The main scale $\lambda$ defines the perturbed objective and appears in the final policy-gradient estimate. The auxiliary scale $\eta$ is used to estimate the population-flow sensitivity $D_t^\theta=\nabla_\theta\mu_t^\theta$. When the exact sensitivity is available, the resulting oracle estimator is unbiased for $\nabla_\theta J^\lambda(\theta)$ by Proposition~\ref{prop:oracle-mse-decomposition}. In the model-free implementation, $D_t^\theta$ is replaced by the recursive estimator $\widehat D_t^{\eta,n,\theta}$.

\begin{algorithm}[H]
\caption{Discrete fixed-scale transport policy gradient}
\label{alg:discrete-transport}
\begin{algorithmic}[1]
\Require policy parameter $\theta$, scales $\lambda,\eta$, main batch size $B$, auxiliary batch size $n$, horizon $T$
\State Compute or estimate the population flow $(\mu_t^\theta)_{t=0}^T$
\State Initialize $\widehat D_0^{\eta,n,\theta}=0$
\State Simulate $n$ auxiliary trajectories with perturbed laws $(1-\eta)\mu_s^\theta+\eta q_s$
\For{$t=1,\ldots,T$}
    \State Estimate $\widehat D_t^{\eta,n,\theta}$ from the shared auxiliary trajectories using the triangular likelihood-ratio recursion
\EndFor
\For{$b=1,\ldots,B$}
    \State Simulate one perturbed trajectory with laws $M_t^{\lambda,\theta,(b)}=(1-\lambda)\mu_t^\theta+\lambda q_t^{(b)}$
    \State Compute the return $\Rc_\theta^{\lambda,(b)}$ and the policy score $\Sc_\mathrm{pol}^{\lambda,\theta,(b)}$
    \State Compute the perturbation score with $\widehat D_t^{\eta,n,\theta}$:
    \begin{align}
        \widehat \Sc_\mathrm{pert}^{\lambda,\eta,n,\theta,(b)}
        =
        -\frac{1-\lambda}{\lambda}
        \sum_{t=0}^T
        \sum_{k=1}^{N-1}
        H_k(q_t^{(b)})
        \widehat D_t^{\eta,n,\theta}(k)
    \end{align}
\EndFor
\State Form
\begin{align}
    \widehat G_{\lambda,\eta,B,n}(\theta)
    =
    \frac1B
    \sum_{b=1}^B
    \left(
        \Rc_\theta^{\lambda,(b)}-b_0
    \right)
    \left(
        \Sc_\mathrm{pol}^{\lambda,\theta,(b)}
        +
        \widehat \Sc_\mathrm{pert}^{\lambda,\eta,n,\theta,(b)}
    \right)
\end{align}
\State Update $\theta$ with Adam in the direction $\widehat G_{\lambda,\eta,B,n}(\theta)$
\end{algorithmic}
\end{algorithm}

The continuous-state transport estimator has the same structure. The difference is that the population law is represented through finite-dimensional features or moments, and the simplex perturbation is replaced by an environment-specific perturbation map. If $Y_t^\theta$ denotes the chosen law feature, the perturbed feature has the generic affine form $Y_t^{\lambda,\theta} = \Phi_\lambda(Y_t^\theta,\xi_t)$,
where $\xi_t$ is exogenous noise. The environment supplies both $\Phi_\lambda$ and the corresponding transport score. The auxiliary stage estimates $\nabla_\theta Y_t^\theta$, and the main stage combines the policy score with the transport score exactly as in Algorithm~\ref{alg:discrete-transport}.

The statement is abstract because the feature $Y_t^\theta$ is whatever finite-dimensional chart the environment exposes. In the two continuous-state benchmarks of Chapter~\ref{chapter:experiments} it is the population mean, the only functional of the law that the coefficients read: $\Phi_\lambda(y,\xi)=(1+\lambda\zeta)y+\lambda\beta$ with $\zeta,\beta$ i.i.d.\ Gaussian, the score is the corresponding Gaussian score of Corollary~\ref{cor:gmm-gaussian-score}, and $\nabla_\theta Y_t^\theta$ is estimated by the auxiliary recursion applied to the closed moment recursion of the benchmark. The exact moment recursion is used to propagate the population during training in the exact-flow runs and is replaced by a particle estimate in the particle-flow runs; in neither case is it differentiated, which is what keeps the estimator model-free in the sense of Section~\ref{sec:canonical-construction}.

\section{Adaptive transport}
\label{sec:adaptive-transport}

Fixed-scale transport requires choosing $\lambda$ and $\eta$ before training. The adaptive version keeps the same gradient estimator but periodically updates these scales from local bias-variance diagnostics. The idea is to compare gradient estimates at the current scales and at contracted scales. If reducing a scale changes the estimated gradient more than expected from Monte Carlo noise, the current scale is interpreted as too biased and is decreased; otherwise it can remain large enough to avoid excessive score variance.

Let $c_\lambda,c_\eta\in(0,1)$ be contraction factors. At an adaptive checkpoint, four gradient estimates are computed using the scale pairs $(\lambda,\eta), \qquad (c_\lambda\lambda,\eta), \qquad (\lambda,c_\eta\eta), \qquad (c_\lambda\lambda,c_\eta\eta)$.
For each replication $r=1,\ldots,R$, denote the corresponding gradient estimates by $g_{++}^{(r)},\qquad g_{-+}^{(r)},\qquad g_{+-}^{(r)},\qquad g_{--}^{(r)}$,
where the first sign refers to the main scale $\lambda$ and the second to the auxiliary scale $\eta$; the sign $+$ means current scale and the sign $-$ means contracted scale. The effect of contracting only $\lambda$ is estimated by averaging over the two auxiliary scales,
\begin{align}
    \Delta_\lambda^{(r)}
    =
    \frac12
    \left[
        \left(g_{++}^{(r)}-g_{-+}^{(r)}\right)
        +
        \left(g_{+-}^{(r)}-g_{--}^{(r)}\right)
    \right],
\end{align}
and similarly the effect of contracting only $\eta$ is
\begin{align}
    \Delta_\eta^{(r)}
    =
    \frac12
    \left[
        \left(g_{++}^{(r)}-g_{+-}^{(r)}\right)
        +
        \left(g_{-+}^{(r)}-g_{--}^{(r)}\right)
    \right].
\end{align}
For a collection of vectors $(U^{(r)})_{r=1}^R$, write
\begin{align}
    \overline U
    =
    \frac1R\sum_{r=1}^R U^{(r)},
    \qquad
    \widehat{\operatorname{trCov}}(U)
    =
    \frac1{R-1}
    \sum_{r=1}^R
    \left\|
        U^{(r)}-\overline U
    \right\|^2,
\end{align}
with $\widehat{\operatorname{trCov}}(U)=0$ if $R=1$. The raw squared discrepancies are corrected for Monte Carlo noise by subtracting the variance of the empirical mean:
\begin{align}
    \widehat d_\lambda
    =
    \left(
        \left\|\overline\Delta_\lambda\right\|^2
        -
        \frac1R
        \widehat{\operatorname{trCov}}(\Delta_\lambda)
    \right)_+,
    \qquad
    \widehat d_\eta
    =
    \left(
        \left\|\overline\Delta_\eta\right\|^2
        -
        \frac1R
        \widehat{\operatorname{trCov}}(\Delta_\eta)
    \right)_+.
\end{align}
If the leading perturbation errors scale like $\lambda^{p_\lambda}$ and $\eta^{p_\eta}$, the bias proxies used by the controller are
\begin{align}
    \widehat b_\lambda
    =
    \frac{\widehat d_\lambda}
    {(1-c_\lambda^{p_\lambda})^2},
    \qquad
    \widehat b_\eta
    =
    \frac{\widehat d_\eta}
    {(1-c_\eta^{p_\eta})^2}.
\end{align}
The variance proxy is the empirical trace covariance of the current-scale estimator, $\widehat v = \widehat{\operatorname{trCov}}(g_{++})$.
The controller then uses the log-ratio signals
\begin{align}
    z_\lambda
    =
    \log
    \left(
        \frac{\widehat b_\lambda+\delta}
        {\tau_\lambda \widehat v+\delta}
    \right),
    \qquad
    z_\eta
    =
    \log
    \left(
        \frac{\widehat b_\eta+\delta}
        {\tau_\eta \widehat v+\delta}
    \right),
\end{align}
where $\tau_\lambda,\tau_\eta>0$ are target bias-to-variance ratios and $\delta>0$ prevents numerical division by zero. Positive values indicate that the estimated perturbation bias is too large relative to variance.

\begin{algorithm}[H]
\caption{Adaptive transport}
\label{alg:adaptive-transport}
\begin{algorithmic}[1]
\Require initial policy $\theta$, initial scales $\lambda,\eta$, contraction factors $c_\lambda,c_\eta$, checkpoint interval
\For{training iterations $m=0,1,\ldots$}
    \State Compute a fixed-scale transport gradient estimate using the current $\lambda,\eta$
    \State Update $\theta$ with Adam
    \If{$m$ is an adaptive checkpoint}
        \State Estimate $g_{++}$, $g_{-+}$, $g_{+-}$, and $g_{--}$ over independent replications
        \State Form $\Delta_\lambda$, $\Delta_\eta$, the corrected discrepancies $\widehat d_\lambda,\widehat d_\eta$, and the variance proxy $\widehat v$
        \State Set $\widehat b_\lambda=\widehat d_\lambda/(1-c_\lambda^{p_\lambda})^2$ and $\widehat b_\eta=\widehat d_\eta/(1-c_\eta^{p_\eta})^2$
        \State Compute $z_\lambda$ and $z_\eta$ from the bias-to-variance ratios
        \State Update transformed scale parameters with an Adam-style controller
        \State Clamp $\lambda$ and $\eta$ to their admissible intervals
    \EndIf
\EndFor
\end{algorithmic}
\end{algorithm}

The adaptive controller is deliberately separated from the policy optimizer. The policy update uses the current transport estimator, while the controller only changes future perturbation scales. Directional checks can also be included: if the current-scale and contracted-scale gradients point in incompatible directions, the corresponding scale is forced downward even when the scalar bias proxy is small.

\section{Computational cost}

The relevant cost unit is one simulated transition. For a horizon $T$ and a main batch of size $B$, REINFORCE has per-update cost $C_{\mathrm{RF}} = BT$.
For finite-state MF-REINFORCE with $n$ auxiliary trajectories and shared state-gradient reuse, the auxiliary procedure requires two perturbed forward simulations for each pair of time indices, yielding $C_{\mathrm{MF}} = BT + nT(T+1)$. Algorithm~1 of \cite{meunier2026model} instead repeats that auxiliary pass inside the loop over the main batch, at a cost of order $nBT^2$; Section~\ref{subsec:experimentation-protocol} explains why the two differ and which is used to match budgets.
The recursive transport estimator uses one auxiliary forward pass and one main forward pass, hence $C_{\mathrm{tr}} = T(B+n)$.
Adaptive transport adds the amortized cost of its checkpoints. If a checkpoint is performed every $K$ iterations and uses $R$ independent replications of the contracted-scale diagnostic, its average per-update cost is written abstractly as $C_{\mathrm{ad}} = C_{\mathrm{tr}} + \frac1K C_{\mathrm{diag}}(R,T)$.

These expressions are used only to account for simulator effort. They do not include common optimizer overhead, tensor operations, logging, or deterministic validation. Their main implication is structural: MF-REINFORCE has quadratic horizon dependence in its auxiliary estimator, while transport keeps linear horizon dependence through the recursive sensitivity construction.
